\documentclass[11pt]{article}
\usepackage[utf8]{inputenc}
\usepackage{amsmath, amssymb, amsthm}
\usepackage{geometry}
\geometry{a4paper, margin=1in}
\usepackage{hyperref}
\usepackage{cite}
\usepackage{graphicx}

\title{Approximation Algorithms and LP Relaxations for Vehicle Routing: \\ Makespan, Latency, and Capacity Constraints}
\author{AI Researcher}
\date{\today}

\begin{document}

\maketitle

\begin{abstract}
This document provides a comprehensive review of mathematical formulations, Linear Programming (LP) relaxations, and approximation algorithms for a class of vehicle routing problems. We consider routing $n$ trucks from a single depot to visit $K$ locations in a metric space. The primary objective is to minimize the arrival time of the last truck to return to the depot (makespan). We also explore extensions involving capacity constraints (where each location demands a certain amount of goods and each truck has a finite capacity) and alternative objectives, such as minimizing the weighted sum of arrival times (latency) at the locations, as well as bi-objective formulations combining makespan and latency.
\end{abstract}

\section{Introduction}

The problem of routing a fleet of vehicles to serve a set of geographically dispersed customers is fundamental in operations research and theoretical computer science. Given a set of $K$ locations (or clients) $V = \{v_1, v_2, \dots, v_K\}$ and a depot $r$ in a metric space with distance function $d(\cdot, \cdot)$, we are tasked with finding a set of tours for $n$ trucks. All trucks start and end at the depot. 

The core version of this problem aims to minimize the \emph{makespan}, defined as the time when the last truck returns to the depot. This problem is known as the Min-Max Multiple Traveling Salesperson Problem (min-max mTSP). It is strictly NP-hard, as it generalizes the classical Traveling Salesperson Problem (TSP) when $n=1$.

\section{Minimizing Makespan: The Min-Max mTSP}

\subsection{Mathematical Formulation and LP Relaxation}
Let $x_{ij}^k \in \{0,1\}$ be a decision variable indicating whether truck $k$ traverses the edge $(i,j)$. Let $y_i^k \in \{0,1\}$ indicate if truck $k$ visits location $i$. 
To minimize the maximum return time $T$, we can formulate a Mixed Integer Linear Program (MILP):

\begin{align}
\min \quad & T \nonumber \\
\text{s.t.} \quad & \sum_{k=1}^n y_i^k = 1, \quad \forall i \in V \label{eq:visit} \\
& \sum_{j} x_{ij}^k = y_i^k, \quad \forall i \in V \cup \{r\}, \forall k \label{eq:degree1} \\
& \sum_{j} x_{ji}^k = y_i^k, \quad \forall i \in V \cup \{r\}, \forall k \label{eq:degree2} \\
& \sum_{i,j \in S} x_{ij}^k \le \sum_{i \in S \setminus \{v\}} y_i^k, \quad \forall S \subseteq V, |S| \ge 2, \forall v \in S, \forall k \label{eq:subtour} \\
& \sum_{i,j} d(i,j) x_{ij}^k \le T, \quad \forall k \label{eq:makespan} \\
& x_{ij}^k \in \{0,1\}, \ y_i^k \in \{0,1\}. \nonumber
\end{seq}
Constraints \eqref{eq:visit} ensure every location is visited exactly once. Constraints \eqref{eq:degree1} and \eqref{eq:degree2} enforce flow conservation. Constraints \eqref{eq:subtour} are the subtour elimination constraints. Constraints \eqref{eq:makespan} bound the total distance traveled by each truck by the makespan $T$.

Relaxing the integrality constraint $x_{ij}^k \in \{0,1\}$ to $0 \le x_{ij}^k \le 1$ yields the standard LP relaxation. However, this relaxation has a large integrality gap for min-max problems. A more powerful relaxation involves the \emph{configuration LP}, which uses a variable for each valid bounded-length tour and has an exponentially large number of variables, solvable via column generation.

\subsection{Approximation Algorithms}
For the min-max mTSP, Even et al.\ (2004) and Franks et al.\ established approximation frameworks based on tree covers. A common approach yields a $3$-approximation (or better, e.g., $2.5$ using Christofides' heuristic) by guessing the optimal makespan $T^*$ via binary search. For a guessed $T$, the algorithm constructs a graph of edges with weight at most $T$ and attempts to find a set of $n$ trees rooted at the depot that cover all vertices such that the weight of each tree is bounded by $\alpha T$. Doubling the trees gives the tours. The current best approximations for metric min-max mTSP are close to $2$.

\section{Incorporating Capacity Constraints (CVRP)}

When each location $i$ demands $g_i$ goods and each truck has a maximum capacity $G$, the problem becomes the Capacitated Vehicle Routing Problem (CVRP) with a min-max objective.

\subsection{LP Relaxation}
We extend the previous LP by adding a constraint ensuring that a truck's total demand does not exceed $G$:
\begin{equation}
\sum_{i \in V} g_i y_i^k \le G, \quad \forall k=1, \dots, n.
\end{equation}
Furthermore, valid inequalities such as the \emph{Capacity Cut} constraints are heavily utilized in the LP relaxation:
\begin{equation}
\sum_{i \in S, j \notin S} x_{ij} \ge 2 \lceil \frac{\sum_{i \in S} g_i}{G} \rceil, \quad \forall S \subseteq V.
\end{equation}
The term $\lceil \sum_{i \in S} g_i / G \rceil$ provides a lower bound on the number of trucks needed to serve subset $S$.

\subsection{Approximation Algorithms}
Approximation algorithms for CVRP often use the Iterated Tour Partitioning (ITP) heuristic originally proposed by Haimovich and Rinnooy Kan (1985). This involves:
1. Solving a standard TSP on all nodes.
2. Partitioning the giant tour into feasible sub-tours such that no sub-tour violates the capacity constraint $G$.
For the min-max objective, Bompadre et al.\ (2006) and others have developed approximations that balance the length of the partitioned segments, leading to constant-factor approximation guarantees that depend on the integrality gap of the underlying TSP relaxation and the bin packing problem.

\section{Minimum Latency and Weighted Arrival Time}

An alternative user-centric objective is to minimize the sum of arrival times (or weighted arrival times) at the locations. This is known as the Minimum Latency Problem (MLP) or the Traveling Repairman Problem (TRP). With $n$ vehicles, it is the $k$-TRP.

\subsection{LP Relaxation}
Formulating MLP mathematically requires capturing the \emph{arrival time} at each node. Let $t_i$ be the arrival time at node $i$. We use a time-indexed or flow-based LP relaxation. A common flow formulation models the number of vertices left to visit. For the multi-vehicle case, configuration LPs where variables represent paths (with their respective latencies) are often used to achieve bounded integrality gaps.

\subsection{Approximation Algorithms}
The first constant-factor approximation for a single vehicle was given by Blum et al.\ (1994). For multiple vehicles ($n$-TRP), Fakcharoenphol, Harrelson, and Rao (2003) and later Chekuri, Korula, and Salavatipour (2012) provided frameworks that yield constant-factor approximations. These algorithms typically involve guessing the locations of the vehicles at exponentially increasing time intervals and solving Prize-Collecting Steiner Tree (PCST) or $k$-MST problems to connect unvisited nodes, thereby controlling the latency accumulation.

\section{Bi-Objective Routing: Makespan and Latency}

In practice, a dispatcher wants to balance the system's efficiency (makespan) with customer satisfaction (latency). 

\subsection{Formulation}
We aim to minimize a convex combination:
\begin{equation}
\min \quad \lambda \cdot (\text{Makespan}) + (1-\lambda) \cdot (\text{Total Weighted Latency})
\end{equation}
for some $\lambda \in [0,1]$. Note that minimizing makespan generally creates sparse, long tours, while minimizing latency encourages dense, short "stars" radiating from the depot.

\subsection{Approximation Algorithms for Bi-Objective}
Finding a single solution that simultaneously approximates both objectives to within a constant factor is challenging because the optimal solutions for the two objectives can be diametrically opposed. However, approximation algorithms exist to approximate the \emph{Pareto frontier}. 
A common technique is the \emph{Bicriteria Approximation}: An algorithm is an $(\alpha, \beta)$-approximation if it produces a solution where the makespan is at most $\alpha$ times the optimal makespan, and the latency is at most $\beta$ times the optimal latency. Constant factor bicriteria approximations for related routing problems (like VRP) leverage modified tree-covers. By optimizing latency subject to a hard constraint on the makespan, rounding the corresponding LP relaxation using lagrangian relaxation techniques yields solutions balancing the two metrics.

\section{Conclusion}

Routing $n$ trucks to service $K$ locations involves complex trade-offs between computational tractability and model fidelity. The LP relaxations for makespan (min-max mTSP) and capacity (CVRP) rely heavily on tree constraints and capacity cuts. Transitioning to latency (MLP) shifts the focus to flow and time-indexed constraints. Approximating a weighted combination of these objectives requires bicriteria techniques that balance the egalitarian nature of makespan against the utilitarian nature of latency. Future work could incorporate stochastic demands or time windows, further enriching the LP structures.

\end{document}
